\begin{question}{29}{
    Instructeurs van de KMA voeren een evaluatie uit om te bepalen of een nieuwe Virtual Reality (VR) schietsimulator net zo effectief is als de traditionele \emph{live-fire} schietbaan voor Close Quarters Battle (CQB) training.
    Hieronder zijn de beoordelingsresultaten van $400$ cadetten gegeven.
    \par\medskip
    \begin{center}
        \renewcommand{\arraystretch}{1}
        \begin{tabular}{cccc}
            \toprule
                & \textbf{Fysieke training} & \textbf{VR-training} & \textbf{Totaal} \\
            \midrule
                Onvoldoende & $24$ & $46$ & $70$  \\
                Voldoende   & $100$ & $132$ & $232$  \\
                Uitmuntend  & $48$ & $50$ & $98$  \\
            \midrule
                \textbf{Totaal} & $172$ & $228$ & $400$ \\
            \bottomrule
        \end{tabular}
    \end{center}
    \par\medskip
    Het doel van de instructeurs is om te bepalen of de keuze voor de trainingsmethode een statistisch significant effect heeft op de schietprestaties.
    Je mag in deze opgave gebruik maken van een significantieniveau $\alpha = 0.05$.
}

    \subquestion{3}{
        Wat voor type hypothesetoets zou geschikt zijn om dit te kunnen toetsen?
        Formuleer de bijbehorende nulhypothese $H_0$ en de alternatieve hypothese $H_1$ voor deze toets.
    }

    \solution{
        Om te bepalen of er een statistisch significant effect is in de trainingsmethode voor de schietprestaties van cadetten,
        dienen we een chikwadraattoets voor onafhankelijkheid uit te voeren. \rubric{1}
        De nulhypothese $H_0$ en de alternatieve hypothesetoets $H_1$ kunnen dan als volgt worden geformuleerd:
        \begin{itemize}
            \item[$H_0$:] de trainingsmethode heeft geen significant effect op de schietprestaties van cadetten. \rubric{1}
            \item[$H_1$:] de trainingsmethode heeft wel een significant effect op de schietprestaties van cadetten. \rubric{1}
        \end{itemize}
    }

    \subquestion{8}{
        Wat is de theoretische toetsingsgrootheid van deze hypothesetoets?
        Bereken de geobserveerde toetsingsgrootheid met behulp van de data in bovenstaande tabel.
    }
    \solution{
        De (theoretische) toetsingsgrootheid voor een chikwadraattoets voor onafhankelijkheid is gelijk aan
        \[
            X^2 = \sum_{i,j} \frac{(O_{ij}-E_{ij})^2}{E_{ij}}.\rubric{1}
        \]    
        Voordat we de toetsingsgrootheid kunnen berekenen, moeten we dus de verwachte frequenties $E_{ij}$ bepalen.
        Dit doen we aan de hand van de volgende formule:
        \[
            E_{ij} = \frac{(\text{rijtotaal}_i) \cdot (\text{kolomtotaal}_j)}{\text{totaal}} \rubric{1}
        \]
        We berekenen de verwachte frequenties voor elke cel in de tabel:
        \par\medskip
        \begin{center}
            \begin{tabular}{c|cc|c}
                \toprule
                    & Fysieke training & VR-training & Totaal \\
                \midrule
                    Onvoldoende & $30.1$  & $39.9$   & $70$ \\ 
                    Voldoende   & $99.76$ & $132.24$ & $232$ \\ 
                    Uitmuntend  & $42.14$ & $55.86$  & $98$ \\ 
                \midrule
                    Totaal & $172$ & $228$ & $400$ \rubric{3} \\
                \bottomrule
            \end{tabular}
        \end{center}
        \par\medskip
        Vervolgens berekenen we de geobserveerde toetsingsgrootheid $\chi^2$ als volgt:
        \begin{align*}
            \chi^2 &= \frac{(o_{1, 1} - e_{1,1})^2}{e_{1, 1}} + \frac{(o_{1, 2} - e_{1,2})^2}{e_{1, 2}} + \ldots + \frac{(o_{3, 2} - e_{3,2})^2}{e_{3, 2}} \\
                   &= \frac{(24 - 30.1)^2}{30.1} + \frac{(46 - 39.9)^2}{39.9} + \ldots + \frac{(50 - 55.86)^2}{55.86} \\
                   &\approx 3.5994. \rubric{3}
        \end{align*}
    }

    \subquestion{5}{
        Wat is de kansverdeling van de theoretische toetsingsgrootheid onder $H_0$? 
        Bereken met deze kansverdeling de $p$-waarde behorende bij de geobserveerde toetsingsgrootheid.
    }
    \solution{
        De theoretische toetsingsgrootheid $X^2$ volgt onder de nulhypothese een chikwadraatverdeling met $\df = (\#\text{rijen}-1)\cdot(\#\text{kolommen}-1) = (3-1)\cdot(2-1) = 2$ vrijheidsgraden. \rubric{2}
        De $p$-waarde is de kans dat de theoretische toetsingsgrootheid $X^2$ een \enquote{nog extremere} uitkomst aanneemt dan de geobserveerde toetsingsgrootheid $\chi^2$. \rubric{1}
        
        Aangezien een chikwadraattoets altijd rechtszijdig is, is de $p$-waarde gelijk aan
        \begin{align*}
            p = P(X^2 > 3.5994) &= \chi^2\text{cdf}(\lb=3.5994, \ub=\GRinfinity, \df=2) \\
                                &\approx 0.1653. \rubric{2}
        \end{align*}

        \begin{center}
            \resizebox{0.9\textwidth}{!}{
                \includegraphics{exam20260724-chi2-pvalue.png} 
            }
        \end{center}
    }

    \subquestion{5}{
        Bepaal de toetsuitslag op basis van de berekende $p$-waarde en vertaal deze terug naar de originele probleemcontext.
        Geef een verklaring voor deze toetsuitslag aan de hand van de gegeven tabel van schietprestaties.
    }
    \solution{
        Aangezien de $p$-waarde $p \approx 0.1653$ strikt groter is dan het significantieniveau $\alpha = 0.05$,
        besluiten we om de nulhypothese $H_0$ niet te verwerpen. \rubric{1}

        Er is onvoldoende reden om aan te nemen dat het type trainingsmethode significant effect heeft op de schietprestaties van cadetten. \rubric{1}
        
        Als we kijken naar de tabellen van geobserveerde en verwachte frequenties,
        dan zien we dat voor veel van de cellen in de tabellen de geobserveerde en verwachte frequenties $o_{ij}$ en $e_{ij}$ niet heel ver uit elkaar liggen.\rubric{2}
        Op basis van deze steekproefdata is er onvoldoende reden om aan te nemen dat het type trainingsmethode significant effect heeft op de schietprestaties van cadetten. \rubric{1}
    }
    
    \subquestion{8}{
        De instructeurs beweren dat de kans dat een willekeurige cadet die de VR-training volgt een voldoende of hoger haalt, minstens \SI{85}{\percent} is.
        Bereken een \SI{95}{\percent}-betrouwbaarheidsinterval voor deze succeskans $p$ met behulp van de Clopper-Pearson methode.
        Wat kunnen we concluderen over de claim van de instructeurs?
    }

    \solution{
        Als we kijken naar het gedeelte van de steekproef dat gaat over cadetten die de VR-training volgen,
        zien we dat $132 + 50 = 182$ van de $228$ cadetten (steekproefschatting!) een voldoende heeft gehaald. \rubric{1}
        Merk op dat $182$ een uitkomst is van de variabele $X$ die het aantal voldoende telt uit $228$ pogingen, oftewel 
        $X$ is binomiaal verdeeld met $n = 228$ en een onbekende $p = \text{ ?}$. \rubric{1}
        Op basis van deze aantallen gaan we een betrouwbaarheidsinterval bepalen voor de kans $p$ dat een willekeurige cadet een voldoende
        na het volgen van de VR-training (populatieparameter!).
        Omdat de gewenste betrouwbaarheid \SI{95}{\percent} is, geldt dat $\alpha = 0.05$.

        De Clopper-Pearson methode werkt als volgt:
        \begin{enumerate}
            \item Bepaal de minimale succeskans $p_1$ waarvoor de uitkomst $k = 182$ niet extreem groot is, oftewel $P(X \geq 182) = \alpha/2 \geq 0.025$.
            Voer hiervoor in de \enquote{numerieke solver} van de grafische rekenmachine in:
            \begin{align*}
                E_1 &= 1 - \binomcdf(n=228, p_1=X, k=181) \\
                E_2 &= 0.025 
            \end{align*}
            De \enquote{numerieke solver} geeft een waarde van $p_1 \approx 0.7402$.
            \item Bepaal de maximale succeskans $p_2$ waarvoor de uitkomst $k = 182$ niet extreem klein is, oftewel $P(X \leq 182) = \alpha/2 \geq 0.025$.
            Voer hiervoor in de \enquote{numerieke solver} van de grafische rekenmachine in:
            \begin{align*}
                E_1 &= \binomcdf(n=228, p_1=X, k=182) \\
                E_2 &= 0.025 
            \end{align*}
            De \enquote{numerieke solver} geeft een waarde van $p_2 \approx 0.8483$.
        \end{enumerate}
        We vinden het Clopper-Pearson interval door de twee gevonden waarden als grenzen te nemen,
        oftewel het \SI{95}{\percent}-betrouwbaarheidsinterval voor de kans $p$ dat een willekeurige cadet een voldoende haalt is gelijk aan $[0.7402, 0.8483]$. \rubric{4}

        We zien dat de geclaimde succeskans $p = 0.85$ niet in het Clopper-Pearson interval ligt. \rubric{1}
        We kunnen dus met \SI{95}{\percent} betrouwbaarheid zeggen dat de kans op een voldoende significant lager is dan \SI{85}{\percent}.\rubric{1}
    }
\end{question}